\section{Ray tracing concepts in the project}
\label{sec:rt}
\here{10}
\concept{A ray $\vect r(t)=\vect o+t\vect d$ is intersected with the scene. Rays answer three questions here: what is
under the cursor (picking), is a point lit (shadow rays), and what does a mirror or glass show (reflection and
refraction rays, Whitted \cite{whitted1980}). The program includes a CPU/GPU ray tracer over the same scene with a BVH.}
\begin{equation}
\text{sphere: } t^2(\vect d\cdot\vect d)+2t\,\vect d\cdot(\vect o-\vect c)+\lVert\vect o-\vect c\rVert^2-r^2=0;\quad
\text{plane: } t=\frac{(\vect p_0-\vect o)\cdot\vect n}{\vect d\cdot\vect n}
\label{eq:rays}
\end{equation}
\begin{equation}
\vect r=\vect d-2(\vect d\cdot\vect n)\vect n,\qquad
\vect t=\eta\,\vect d+\big(\eta\cos\theta_i-\sqrt{1-\eta^2(1-\cos^2\theta_i)}\big)\vect n,\ \eta=\tfrac{n_1}{n_2}
\label{eq:reflref}
\end{equation}

\begin{figure}[H]
\centering
\raw{rt/raytrace_modes}{\linewidth}
\caption{The same view in the five render modes (\key{;}): raster, hybrid (raster plus ray-traced reflections),
Whitted, path tracing and light tracing. Notice the shiny surfaces: the raster image can only reflect what is on screen, Whitted traces exact
mirror reflections and shadow rays, and the path tracer adds soft indirect light and soft shadows.}
\label{fig:rtmodes}
\end{figure}

\begin{figure}[H]
\centering
\begin{minipage}{0.5\linewidth}\centering\raw{rt/raytrace_tree}{\linewidth}\end{minipage}\hfill
\begin{minipage}{0.48\linewidth}\centering\wb{Ray_Tracing}{\linewidth}\end{minipage}
\caption{Left: the ray debugger (\key{'}, click a pixel) draws the ray tree of one pixel: primary ray, shadow rays to
each light, reflection and refraction children. Right: workbook sheet \code{Ray\_Tracing}: ray--sphere and
ray--plane intersections, reflection, Snell refraction and the contribution of each recursion level.}
\label{fig:rttree}
\end{figure}

\textbf{Uses in the game.} Mouse picking casts a ray through the clicked pixel (middle click jumps into that object);
the interaction prompt tests a segment against wall boxes for line of sight; the UAV searchlight's beam length is a
ray marched to the terrain and refined by bisection; a third-person camera arm is a sphere cast against the walls.
\where{src/raytrace/ (Intersect, Bvh, RayTracer), src/rendering/RayTracerGL.cpp, src/scene/Interact.cpp (segmentClear), src/scene/Scene.cpp (beam length)}
\critique{The ray-traced modes are teaching modes: at 1080p the path tracer runs at reduced internal resolution
(0.5) and accumulates over frames, so it is not used in play. The raster renderer approximates reflections with
screen-space reflections (Section~\ref{sec:rtadv}).}
